% =========================
%  Insulation Coordination
% =========================
\intermezzo{Arrester}{}

\begin{frame}{Overvoltages in AC grids}
  \centering
  \topoimg[width=\linewidth,height=0.72\textheight,keepaspectratio]{\icfig{m08_image7}}\\[0.3em]
  {\color{atGray}\footnotesize ``Metal-Oxide Surge Arresters in High-Voltage
  Power Systems'', Volker Hinrichsen.}
\end{frame}

\begin{frame}{Arrester design}
  \begin{columns}[c,onlytextwidth]
    \begin{column}{0.56\textwidth}
      \begin{itemize}
        \setlength{\itemsep}{0.4em}
        \item Arresters used today are metal-oxide (MO) arresters without gaps.
        \item Made from MO resistor disks connected in series.
        \item The number of series-connected resistors defines the arrester
              characteristics.
        \item MO resistors have an extremely non-linear $U-I$ characteristic.
        \item Resistors are placed in a housing for protection against
              environmental conditions.
      \end{itemize}
    \end{column}
    \begin{column}{0.40\textwidth}
      \centering
      \topoimg[width=\linewidth,height=0.68\textheight,keepaspectratio]{\icfig{m08_image8}}
      {\color{atGray}\footnotesize ``Metal-Oxide Surge Arresters in High-Voltage Power Systems``, Volker Hinrichsen}
    \end{column}
  \end{columns}
\end{frame}

\begin{frame}{U--I arrester characteristic}
  \begin{columns}[c,onlytextwidth]
    \begin{column}{0.54\textwidth}
      \begin{itemize}
        \setlength{\itemsep}{0.35em}
        \item $U_S$: maximum system voltage (phase-phase AC).
        \item $U_c$: maximum continuous operating voltage (MCOV).
        \item For DC applications the crest value of continuous operating
              voltage (CCOV) is used instead of MCOV.
        \item Lightning impulse current ($I_n$): $8/20 \mu s$
              (fast-front).
        \item Switching impulse current: $30/60 \mu s$ (slow-front).
      \end{itemize}
    \end{column}
    \begin{column}{0.42\textwidth}
      \centering
      \topoimg[width=\linewidth,height=0.68\textheight,keepaspectratio]{\icfig{m08_image9}}
    \end{column}
  \end{columns}
\end{frame}

\begin{frame}{Arrester protection capability --- connection}
  \begin{columns}[c,onlytextwidth]
    \begin{column}{0.56\textwidth}
      \begin{itemize}
        \setlength{\itemsep}{0.4em}
        \item Protection capability is based on providing a low-impedance
              path towards earth.
        \item That impedance is defined by the arrester itself and by its
              connection to the overhead line and to earth.
        \item Either accurate calculations or safety margins are applied to
              account for those effects.
      \end{itemize}
    \end{column}
    \begin{column}{0.40\textwidth}
      \centering
      \topoimg[width=\linewidth,height=0.66\textheight,keepaspectratio]{\icfig{m08_image10}}
    \end{column}
  \end{columns}
\end{frame}

\begin{frame}{Arrester protection capability - travelling wave}
  \begin{columns}[c,onlytextwidth]
    \begin{column}{0.56\textwidth}
      \begin{itemize}
        \setlength{\itemsep}{0.4em}
        \item The surge current will enter a ``dead end'' (e.g. transformer):
              even with an arrester installed, the travelling wave enters the
              transformer connection and causes voltage stress.
        \item The resulting stress is an overlap of the incoming surge, its
              reflection and the arrester protection rating.
        \item Equipment voltage stress therefore depends not only on the
              arrester protective level but also on the distance between
              arrester and protected device.
      \end{itemize}
    \end{column}
    \begin{column}{0.40\textwidth}
      \centering
      \topoimg[width=\linewidth,height=0.66\textheight,keepaspectratio]{\icfig{m08_image11}}
      {\color{atGray}\footnotesize ``Metal-Oxide Surge Arresters in High-Voltage Power Systems``, Volker Hinrichsen}
    \end{column}
  \end{columns}
\end{frame}

\begin{frame}{Energy handling capability}
  \begin{columns}[c,onlytextwidth]
    \begin{column}{0.56\textwidth}
      \begin{itemize}
        \setlength{\itemsep}{0.4em}
        \item \textbf{Simple impulse handling} --- energy injected in only a
              few \textmu s must not cause mechanical damage.
        \item \textbf{Thermal energy handling} --- maximum energy injected
              during a transient at which the arrester can still cool back to
              its nominal operating temperature.
        \item Not precisely defined in the standards, so for VSC-HVDC the
              energy loading of the arrester is an important design point.
      \end{itemize}
    \end{column}
    \begin{column}{0.40\textwidth}
      \centering
      \topoimg[width=\linewidth,height=0.66\textheight,keepaspectratio]{\icfig{m08_image12}}
      {\color{atGray}\footnotesize ``Metal-Oxide Surge Arresters in High-Voltage Power Systems``, Volker Hinrichsen}
    \end{column}
  \end{columns}
\end{frame}

\begin{frame}{Arrester configuration in VSC-HVDC systems}
  \begin{columns}[c,onlytextwidth]
    \begin{column}{0.48\textwidth}
      \centering
      \topoimg[width=\linewidth,height=0.66\textheight,keepaspectratio]{\icfig{m08_image15}}
    \end{column}
    \begin{column}{0.48\textwidth}
      \centering
      \topoimg[width=\linewidth,height=0.66\textheight,keepaspectratio]{\icfig{m08_image16}}
    \end{column}
  \end{columns}
  \vspace{0.3em}
  \centering
  {\color{atGray}\footnotesize IEC 60071-11:2022.}
\end{frame}
